Le passage de l’hexade à l’octade est une transformation de la structure d’incidence.\par

Lors du passage de \(6\) à \(8\), les paires disponibles, les relations internes et la manière dont l’orientation peut se conserver changent. Cette transition engendre un défaut orienté. La variation cardinale réalise une rotation du mode d’organisation.\par

La même inversion de feuillet s’exprime par\par

\[
C^*\longmapsto -C^*.
\]

Mais des lecteurs différents répondent de manières différentes à cette inversion.\par

La fonctionnelle hexade–octade qui conduit au paramètre de Barbero–Immirzi est impaire : lorsque le feuillet est inversé, elle change de signe. L’orientation géométrique s’inverse.\par

L’entropie bicouche, en revanche, est paire : elle attribue la même entropie à \(C^*\) et à \(-C^*\) lorsque l’on observe la distribution de probabilités. L’information totale peut rester identique même si l’orientation a été inversée.\par

L’espace CKM réalise une troisième réponse. L’inversion induit sur les trois angles une réflexion rationnelle exacte :\par

\[
R_{\rm CKM}
=
M_{\rm CKM}
\operatorname{diag}(1,-1,1)
M_{\rm CKM}^{-1},
\]

con\par

\[
R_{\rm CKM}^2=I,
\qquad
\det R_{\rm CKM}=-1.
\]

Cela signifie que le changement de feuillet est transporté dans l’espace de mélange comme une véritable réflexion. Il existe une direction impaire de rang un et un plan qui reste fixe.\par

Le contre-angle \(C^*\) est précisément la coordonnée qui change. La phase construite à partir de \(A\) reste invariante. La transformation sépare ainsi deux mémoires :\par

\begin{itemize}[leftmargin=*,itemsep=0.35em]
\item une mémoire d’orientation ;
\item une mémoire de phase.
\end{itemize}

Cela offre une lecture très profonde de CKM. Les angles de mélange sont les coordonnées d’un espace sur lequel agit la même involution qui organisait déjà la transition hexade–octade.\par

La rotation \(6\to8\) fournit l’opération d’orientation que CKM réalise dans son propre codomaine ; la causalité physique appartient au transducteur correspondant.\par

Nous avons alors trois publications sœurs du même passage :\par

\begin{enumerate}[leftmargin=*,itemsep=0.65em]
\item \textbf{Publication géométrique.}
Le paramètre de Barbero–Immirzi oriente le quantum d’aire.\par

\item \textbf{Publication informationnelle.}
L’état réduit et son hamiltonien modulaire mesurent quelle information de cette orientation demeure au bord.\par

\item \textbf{Publication de mélange.}
La réflexion CKM transforme les coordonnées de mélange et isole leur composante impaire.\par
\end{enumerate}

\(\gamma_{\rm BI}\), l’entropie et les angles CKM demeurent des réalisations distinctes d’une charnière généalogique commune. L’unité réside dans l’antécédent antérieur à la ramification. Ils partagent la charnière, mais chacun en conserve une partie différente.\par

C’est l’une des plus belles expressions de la méthodologie HMT : l’unité fondamentale est retrouvée en remontant les généalogies jusqu’à trouver l’opération commune qui engendre des résultats terminaux différents.\par

La gravité conserve la rotation comme orientation de l’aire.\par
L’information conserve la rotation comme mémoire de l’état.\par
Le mélange conserve la rotation comme réflexion des coordonnées.\par

La même transformation « respire » de trois manières.\par

